%%%%%%%%%%  scan idx 325  |  folio 312  %%%%%%%%%%
(1.7) que \(\Phi\) est
déterminé par sa valeur sur l'élément \((\varphi^{(j)})_{j=1..k}\) de \((\Ghat(K[X_i^{(j)}]_{\substack{i\geq1\\ j=1..k}}))^k\) (resp.\ \ldots)..\ Supposons que \(H\) soit un foncteur tel que \(\Ghat,\Gg,\Ga,1+\widehat{\operatorname{Gr}\,\cdot\,}^{+}\), etc..~; l'élément \(\Phi((\varphi^{(j)})_{j})\) est une série formelle à coefficients dans \(K[X_i^{(j)}]_{i\geq1,\ j=1..k}\)~; le \(n\)-ième coefficient de cette série sera appelé \emph{\(n\)-ième polynôme universel de} \(\Phi\)~; la connaissance de ces polynômes équivaut donc à la connaissance de \(\Phi\). D'après la démonstration de 1.6, ces polynômes sont eux-mêmes déterminés par l'élément \(\Phi(1+X^{(1)}t,\ldots,1+X^{(k)}t)\), le calcul s'effectuant selon la méthode de la démonstration : on calcule les éléments \(\Phi(1+X_1^{(1)}t+\ldots+X_n^{(1)}t^n,\ldots,1+X_1^{(k)}t+\ldots+X_n^{(k)}t^n)\), en décomposant \(1+X_1^{(i)}t+\ldots+X_n^{(i)}t^n\) en \(\prod_{j=1}^{n}(1+T_j^{(i)}t)\) et en utilisant l'additivité de \(\Phi\)~; les polynômes ainsi obtenus sont symétriques par rapport aux \(T_j^{(i)}\) (\(i\) fixe et \(j\) variable) et s'expriment en fonction des \(X_j^{(i)}\)~; en recollant grâce à la filtration de \(H\), on détermine ainsi les polynômes universels de \(\Phi\). Si \((\psi^{(1)},\ldots,\psi^{(k)})\in(\Ghat(A))^k\), le calcul de \(\Phi(\psi^{(1)},\ldots,\psi^{(k)})\) se fait en substituant les coefficients \(a_j^{(i)}\) des \(\psi^{(i)}\) aux \(X_j^{(i)}\) dans les polynômes universels.

Supposons qu'il existe un entier \(m\) tel que le coefficient de \(t^n\) dans \(\theta=\Phi(1+X^{(1)}t,\ldots,1+X^{(k)}t)\) soit un polynôme dont les composantes isobares ont un poids variant entre \(n\) et \(mn\)~; alors le \(k\)-ième polynôme universel de \(\Phi\) a des composantes isobares dont le poids varie entre \(k\) et \(km\)~: ceci résulte immédiatement du mode de calcul exposé plus haut, et du théorème des fonctions symétriques.

